Let a route-ranking request be defined by origin $o$, destination $d$, request context $c$, optional preference text $p$, and optional historical records $H_u$ for a user or pseudo-user. The system must generate a candidate set $R = \{r_1,\ldots,r_k\}$ and return an ordering that places routes aligned with observed or inferred preferences near the top.

Each route $r_i$ is represented by an interpretable feature vector:

\begin{equation}
x_i = [\mathrm{distance}, \mathrm{major}, \mathrm{walk}, \mathrm{residential}, \mathrm{service}, \mathrm{turns}, \ldots].
\end{equation}

For OSM-derived pseudo-history evaluation, the held-out record is not treated as a generated candidate. Instead, it is treated as an observed reference route with feature vector $x^\ast$. The oracle generated candidate is:

\begin{equation}
r^\ast = \arg\min_{r_i \in R} \mathrm{dist}(\hat{x}_i, \hat{x}^\ast),
\end{equation}

where $\hat{x}$ denotes min-max normalized features within the evaluation set.

The research question is:

\begin{quote}
Can dynamic profiles built from OSM-derived historical movement signals rank generated OSM route candidates closer to observed pseudo-route behavior than non-personalized baselines?
\end{quote}

This is intentionally narrower than claiming full personalized navigation. The paper evaluates a reranking problem: candidate routes are already generated by OSM graph search, and the proposed profile module changes their order. This distinction matters because a learned search method such as NASR can generate different paths, while our current implementation can only choose among candidates in the generated pool.

\subsection{Evaluation Assumptions}
The evaluation uses two kinds of observed reference routes. In the OSM public-trace setting, references are reconstructed walking-route pseudo-histories from public GPS trackpoints. In the Porto setting, references are reconstructed driving routes from public taxi trajectories. In both cases, the observed reference is imperfect. We therefore report reconstruction diagnostics and oracle upper bounds rather than treating the held-out route as unquestioned ground truth.

\subsection{Baseline Objective}
For each held-out query, a method returns an ordering of generated candidates. A successful method should rank the oracle candidate near the top, where the oracle is either the candidate closest in normalized route-feature space or, for path diagnostics, the candidate with strongest overlap with the reconstructed reference path. The main deployable baselines are random, shortest-distance, profile, and vehicle-profile ranking. Feature and path oracles are diagnostic upper bounds, not usable methods.
